\documentclass[11pt,a4paper]{article}

% ---------- Packages ----------
\usepackage[margin=1in]{geometry}
\usepackage{amsmath, amssymb, amsthm, mathtools}
\usepackage{microtype}
\usepackage{algorithm}
\usepackage{algpseudocode}
\usepackage[colorlinks=true,
            linkcolor=blue!50!black,
            urlcolor=blue!50!black,
            citecolor=blue!50!black]{hyperref}

% ---------- Theorem environments ----------
\theoremstyle{plain}
\newtheorem{theorem}{Theorem}[section]
\newtheorem{proposition}[theorem]{Proposition}
\newtheorem{lemma}[theorem]{Lemma}
\newtheorem{corollary}[theorem]{Corollary}

\theoremstyle{definition}
\newtheorem{definition}[theorem]{Definition}

\theoremstyle{remark}
\newtheorem{remark}[theorem]{Remark}
\newtheorem{example}[theorem]{Example}

% ---------- Math shortcuts ----------
\newcommand{\R}{\mathbb{R}}
\newcommand{\N}{\mathbb{N}}
\newcommand{\E}{\mathbb{E}}
\newcommand{\Var}{\operatorname{Var}}
\newcommand{\F}{\mathcal{F}}
\newcommand{\Prob}{\mathbb{P}}
\newcommand{\Q}{\mathbb{Q}}
\newcommand{\dd}{\mathop{}\!\mathrm{d}}
\newcommand{\eqdef}{\overset{\mathrm{def}}{=}}

% ---------- Document metadata ----------
\title{Monte Carlo Pricing of a European Call:\\
       Euler--Maruyama Discretisation under Geometric Brownian Motion}
\author{Quant Finance Portfolio --- Phase 2, Block 1, Algorithm 2}
\date{\today}

\begin{document}
\maketitle

\begin{abstract}
We specialise the Euler--Maruyama scheme of the previous theoretical
writeup to geometric Brownian motion, derive the resulting recursion,
and apply it to Monte Carlo pricing of the European call. The
algorithm is biased: unlike the exact sampler of Block~1.1, the
estimator converges to $\mathbb{E}[\Phi(\hat S_T)]$, not to
$\mathbb{E}[\Phi(S_T)]$, and the gap is the discretisation bias of
order $h$ predicted by the weak convergence theorem. The principal
contribution of this writeup is the design of the empirical
convergence study that validates the predicted strong order
$\tfrac{1}{2}$ and weak order $1$ on a controlled benchmark, anchored
on the closed-form GBM solution that allows pathwise comparison
between the discretised and the exact dynamics.
\end{abstract}

\tableofcontents

\section{Introduction}
\label{sec:intro}

The previous writeup (\texttt{sde\_discretisation.tex}) developed the
abstract theory of strong and weak convergence for SDE discretisation
schemes, derived the Euler--Maruyama scheme as a first-order
truncation of the It\^o-Taylor expansion, and stated its convergence
orders: strong $\tfrac{1}{2}$ and weak $1$. The present document
specialises that theory to the simplest non-trivial setting in
finance: pricing a European call under geometric Brownian motion,
where the underlying SDE admits a closed-form solution and a direct
comparison between discretised and exact paths is therefore possible.

The setting is degenerate in the sense that no one would actually use
Euler to price a European call under GBM: the Block~1.1 algorithm,
which samples $S_T$ exactly with one normal draw per path, is both
faster (no inner time loop) and exact (no discretisation bias). The
value of the present algorithm is twofold. First, it is the simplest
laboratory in which to verify empirically that the convergence orders
predicted by Theorem~5.1 of the previous writeup are real and
quantitatively correct, on a problem where the exact solution is
known and pathwise comparison is feasible. Second, the Euler scheme
implemented here is the foundation on which all SDE-driven Monte
Carlo work in subsequent blocks rests: when local volatility
(Block~5), or jump-diffusion, or stochastic volatility models are
added, the structure of the pipeline (sample increments, recurse the
state, evaluate payoff at terminal value, average) is unchanged; only
the recursion changes.

The remainder of this document is organised as follows.
Section~\ref{sec:setting} restates the SDE setting and specialises
Euler--Maruyama to GBM. Section~\ref{sec:estimator} defines the
biased Monte Carlo estimator and decomposes its mean square error
into bias and variance contributions.
Section~\ref{sec:validation} designs the empirical convergence study.
Section~\ref{sec:algorithm} states the algorithm formally.
Section~\ref{sec:practical} addresses the practical question of
choosing $h$ at fixed computational budget.

\section{Euler--Maruyama for geometric Brownian motion}
\label{sec:setting}

\subsection{The recursion}

Recall the GBM SDE under the risk-neutral measure $\Q$:
\begin{equation}
\label{eq:sde-gbm}
\dd S_t \;=\; r S_t \dd t + \sigma S_t \dd W_t,
\qquad S_0 > 0.
\end{equation}
With $a(s) = r s$, $b(s) = \sigma s$, the Euler--Maruyama
recursion~(\ref{eq:euler}) of the previous writeup specialises to
\begin{equation}
\label{eq:euler-gbm}
\hat S_{n+1} \;=\; \hat S_n \;+\; r\, \hat S_n\, h
\;+\; \sigma\, \hat S_n\, \Delta W_n
\;=\; \hat S_n\,\big(\, 1 + r h + \sigma\, \Delta W_n \,\big),
\qquad \hat S_0 = S_0,
\end{equation}
with the Brownian increments $\Delta W_n =
W_{(n+1) h} - W_{n h} \sim \mathcal{N}(0, h)$ independent across
$n = 0, 1, \ldots, N - 1$, and $h = T / N$.

The recursion is multiplicative: $\hat S_{n+1}$ is $\hat S_n$
multiplied by the random factor $(1 + r h + \sigma \Delta W_n)$. As a
consequence,
\begin{equation}
\hat S_T \;=\; \hat S_N \;=\;
S_0 \prod_{n=0}^{N-1} \big(\, 1 + r h + \sigma \Delta W_n \,\big).
\end{equation}
This is a discrete-time, geometric-multiplicative process.

\subsection{Failure of positivity}

The exact GBM solution is positive almost surely: $S_t = S_0
\exp((\mu - \tfrac{1}{2}\sigma^2) t + \sigma W_t) > 0$. The Euler
recursion does not preserve this property. Whenever
$1 + r h + \sigma \Delta W_n < 0$, we have $\hat S_{n+1}$ of the
\emph{opposite sign} from $\hat S_n$. This requires
\begin{equation}
\Delta W_n \;<\; -\frac{1 + r h}{\sigma},
\end{equation}
an event that, since $\Delta W_n \sim \mathcal{N}(0, h)$, has
probability
\begin{equation}
\Phi\!\left(-\frac{1 + r h}{\sigma \sqrt{h}}\right).
\end{equation}
For $r = 0.05$, $\sigma = 0.20$, and $h = 1/100$, this probability is
approximately $\Phi(-50.25) \approx 10^{-553}$: astronomically small,
operationally zero. For $h = 1/10$ and the same parameters, the
threshold is $\Phi(-15.96) \approx 10^{-58}$: still operationally
zero. For very large $\sigma$ or very small $N$ the threshold relaxes
but remains far in the tail.

We will not implement any defence against this event in the algorithm
itself: the empirical probability is small enough that no run we will
perform is expected to encounter it. If the simulation grids are
extended to include very coarse $h$ or very high volatility, a
positivity check can be added at little cost. A numerically prudent
alternative is the \emph{log-Euler} scheme, which applies Euler in
$\log S$ and then exponentiates; for GBM this turns out to coincide
with the exact sampler of Block~1.1, so it is not informative as a
test of the Euler convergence rate.

\subsection{Convergence orders}

By Theorem~5.1 of the previous writeup, applied to the GBM
coefficients which satisfy the regularity assumptions, the Euler
recursion~\eqref{eq:euler-gbm} converges to the exact $S_T$ with
\begin{equation}
\label{eq:euler-orders}
\E\big[\,|S_T - \hat S_T|\,\big] \;=\; O(h^{1/2}),
\qquad
|\E[f(S_T)] - \E[f(\hat S_T)]| \;=\; O(h)
\quad \text{for } f \text{ smooth enough.}
\end{equation}
The first is strong order $\tfrac{1}{2}$, the second is weak order
$1$. Both will be verified empirically in
Section~\ref{sec:validation}.

\section{The biased Monte Carlo estimator}
\label{sec:estimator}

\subsection{Definition}

Let $\hat S_T^{(i)}$ for $i = 1, \ldots, M$ denote independent
realisations of the Euler-discretised terminal price, each generated
by the recursion~\eqref{eq:euler-gbm} with its own sequence of
Brownian increments. The Monte Carlo Euler estimator of the
European call price is
\begin{equation}
\label{eq:estimator}
\hat C_{M, h}^{\mathrm{Euler}} \;\eqdef\; \frac{1}{M} \sum_{i=1}^{M}
e^{-rT}\, \big(\hat S_T^{(i)} - K\big)^+ .
\end{equation}
Two sources of error contribute to $\hat C_{M, h}^{\mathrm{Euler}}
- C_{\mathrm{BS}}$: a discretisation bias from replacing $S_T$ with
$\hat S_T$, and a sampling error from finite $M$.

\subsection{Bias--variance decomposition}

Define
\begin{equation}
C_{\mathrm{Euler}}(h) \;\eqdef\; \E\big[\,e^{-rT}\,(\hat S_T - K)^+\,\big],
\end{equation}
the population mean of the discretised payoff at step size $h$. The
estimator decomposes as
\begin{equation}
\hat C_{M, h}^{\mathrm{Euler}} - C_{\mathrm{BS}}
\;=\; \big(\hat C_{M, h}^{\mathrm{Euler}} - C_{\mathrm{Euler}}(h)\big)
\;+\; \big(C_{\mathrm{Euler}}(h) - C_{\mathrm{BS}}\big).
\end{equation}
The first term has zero mean and standard deviation
$\sqrt{\Var(\hat S_T \cdot \text{payoff}) / M}$, by the same central
limit theorem argument as in Block~1.1: it is the statistical noise
controlled by sample size $M$. The second term is deterministic and
does not vanish as $M \to \infty$: it is the discretisation bias
controlled by step size $h$, and by~\eqref{eq:euler-orders} it is
$O(h)$.

The mean square error of the estimator is therefore
\begin{equation}
\label{eq:mse-decomp}
\E\big[\,(\hat C_{M, h}^{\mathrm{Euler}} - C_{\mathrm{BS}})^2\,\big]
\;=\; \frac{\sigma_{\mathrm{Euler}}^{\,2}(h)}{M} \;+\;
\big(C_{\mathrm{Euler}}(h) - C_{\mathrm{BS}}\big)^{\!2}
\;\sim\; \frac{C_1}{M} + C_2\, h^{\,2},
\end{equation}
where $\sigma_{\mathrm{Euler}}^{\,2}(h)$ is the variance of the
Euler-discretised discounted payoff (which converges to the exact
$\Var(Y)$ of Block~1.1 as $h \to 0$, hence the constant approximation
$C_1$), and the second term reflects the squared bias.

This decomposition is the conceptual novelty of Block~1.2 with
respect to Block~1.1, where the bias was identically zero. Choosing
the simulation parameters now requires balancing two effects rather
than one. We address the practical implications in
Section~\ref{sec:practical}.

\section{Empirical validation strategy}
\label{sec:validation}

This section is the heart of the document. We design four tests that
together verify both the theoretical predictions of
Theorem~5.1 of the previous writeup and the implementation
correctness. Two of the tests measure convergence orders directly
(strong and weak); the third anchors the implementation in the
existing exact pricer; the fourth checks a basic invariance.

\subsection{Test 1: empirical strong order}

The strong order is the rate of decay of $\E[|S_T - \hat S_T|]$ in
$h$. To estimate this expectation, we need to evaluate $S_T$ and
$\hat S_T$ \emph{on the same Brownian path}. The trick that makes this
feasible is the closed-form GBM solution: given the simulated
Brownian increments $\Delta W_0, \ldots, \Delta W_{N-1}$ used by the
Euler scheme, the corresponding exact terminal value is
\begin{equation}
\label{eq:exact-from-W}
S_T^{\mathrm{exact}} \;=\; S_0 \exp\!\left(
\big(r - \tfrac{1}{2}\sigma^2\big) T + \sigma\, W_T
\right),
\qquad
W_T = \sum_{n=0}^{N-1} \Delta W_n.
\end{equation}
This $S_T^{\mathrm{exact}}$ is the value of the true SDE solution at
$T$ for the very same realisation of Brownian motion that drove the
Euler simulation. The pathwise difference $|S_T^{\mathrm{exact}} -
\hat S_T|$ is therefore a sample of the strong error.

The protocol is:
\begin{enumerate}
\item Choose a logarithmic grid of step counts $N_1 < N_2 < \ldots <
N_K$, equivalently $h_k = T/N_k$.
\item For each $N_k$, generate $M$ independent Euler paths, each with
its own sequence of $N_k$ Brownian increments.
\item For each path $i$, compute $S_T^{\mathrm{exact}, (i)}$
from~\eqref{eq:exact-from-W} and $\hat S_T^{(i)}$ from the Euler
recursion.
\item Estimate $\mathrm{strong}(h_k) \approx M^{-1} \sum_{i=1}^{M}
|S_T^{\mathrm{exact}, (i)} - \hat S_T^{(i)}|$.
\item Plot $\log(\mathrm{strong}(h_k))$ against $\log(h_k)$.
\end{enumerate}
The expected behaviour is a straight line of slope close to
$\tfrac{1}{2}$. Substantial deviation from this slope indicates
either a bug in the Euler implementation, in the exact comparison, or
a regime where the asymptotic order has not yet been reached (very
small $N$).

\subsection{Test 2: empirical weak order}

The weak order is the rate of decay of $|\E[f(S_T)] - \E[f(\hat S_T)]|$
in $h$, for the call payoff $f(s) = e^{-rT}(s - K)^+$ in our setting.

\begin{enumerate}
\item Choose the same grid of step counts $N_k$.
\item For each $N_k$, generate $M$ Euler paths and form
$\hat C_{M, h_k}^{\mathrm{Euler}}$ from~\eqref{eq:estimator}.
\item Estimate $\mathrm{weak}(h_k) \approx |\hat C_{M, h_k}^{\mathrm{Euler}}
- C_{\mathrm{BS}}|$.
\item Plot $\log(\mathrm{weak}(h_k))$ against $\log(h_k)$, fit a
straight line.
\end{enumerate}
The expected slope is close to $1$. The two convergence studies
together verify the asymmetry of the SDE convergence orders, the
single most important conceptual takeaway of Block~1.2.

\subsection{Calibration of $M$ for the convergence studies}

Tests~1 and~2 are subtle: they estimate quantities that are
themselves expectations, so the Monte Carlo noise $O(M^{-1/2})$ acts
as a floor on what slopes can be measured.

For the strong test, the empirical strong error
$\mathrm{strong}(h)$ has standard error of order
$\sqrt{\Var(|S_T^{\mathrm{exact}} - \hat S_T|) / M}$. As long as this
standard error is small compared with the strong error itself, the
slope is well-defined. For $h \in [10^{-4}, 10^{-1}]$ and our GBM
parameters, $|S_T^{\mathrm{exact}} - \hat S_T|$ has typical scale
$10^{-1}$ to $10^{-2}$ and standard deviation comparable; setting
$M = 10^4$ gives a relative standard error of $10^{-2}$, sufficient
to resolve the $\tfrac{1}{2}$ slope across two orders of magnitude in
$h$.

For the weak test, the empirical weak error
$\mathrm{weak}(h)$ has standard error of order $\sigma_Y / \sqrt{M}$
where $\sigma_Y$ is the standard deviation of the discounted payoff.
For our GBM parameters $\sigma_Y \approx 14$ and the weak error is
$O(h)$ which can be as small as $10^{-3}$. To resolve the slope $1$
we need
$\sigma_Y / \sqrt{M} \ll 10^{-3}$, i.e.\ $M \gg 2 \times 10^8$.
This is impractical. The standard remedy is \emph{common random
numbers between the simulated and the reference computations}: instead
of comparing $\hat C_{M, h}^{\mathrm{Euler}}$ with the closed-form
$C_{\mathrm{BS}}$, we compute both estimators (Euler and exact) on
the same Brownian paths and look at their difference, which has much
smaller variance than either estimator individually.

Specifically, with the same $M$ Brownian paths, evaluate
\begin{equation}
\hat D_{M, h} \;\eqdef\;
\frac{1}{M} \sum_{i=1}^{M} e^{-rT}
\left[ (\hat S_T^{(i)} - K)^+ - (S_T^{\mathrm{exact}, (i)} - K)^+ \right].
\end{equation}
Then $\E[\hat D_{M, h}] = C_{\mathrm{Euler}}(h) - C_{\mathrm{BS}}$,
which is exactly the weak error, and
$\Var(\hat D_{M, h}) = M^{-1} \Var(\Phi(\hat S_T) - \Phi(S_T))$,
which is much smaller than $M^{-1} \Var(\Phi(\hat S_T))$ because the
two payoffs are highly correlated under common Brownian. With $M =
10^5$ this is sufficient to resolve the slope.

\subsection{Test 3: coherence with the exact pricer}

This test is operational rather than theoretical: it confirms that as
$h \to 0$ at fixed $M$, the Euler pricer's estimate approaches the
Block~1.1 exact pricer's estimate (modulo the irreducible $O(M^{-1/2})$
sampling noise). A failure here would signal a bug, not a violation
of the theory.

The cleanest way to run this test under common random numbers is:
generate one batch of $M$ Brownian paths at $N$ steps each, run the
Euler scheme on them, and compare the resulting estimator against the
exact estimator built from $W_T = \sum_n \Delta W_n$ on the same
paths. The two estimators share noise; their difference is at most
the discretisation bias.

\subsection{Test 4: input validation and basic invariances}

A separate, mechanical battery of tests checks that the implementation
rejects bad inputs (negative $S$, $K$, $\sigma$, $T$, non-positive
$N$ or $M$), that the result is monotone in the parameters in the
expected directions (modulo noise), and that the asymptotic limits
$S_0 \to \infty$ and $S_0 \to 0$ are respected, mirroring Test~4 of
Block~1.1.

\section{Algorithmic specification}
\label{sec:algorithm}

The algorithm is summarised below. The pseudocode is parametric in
the random source, following the seed-control discipline established
in Block~0.

\begin{algorithm}[H]
\caption{Euler--Maruyama Monte Carlo for the European call (GBM)}
\label{alg:mc-euler}
\begin{algorithmic}[1]
\Require Spot $S_0$, strike $K$, maturity $T$, rate $r$, volatility
$\sigma$, step count $N$, sample size $M$, confidence level
$1 - \beta$, random generator \texttt{rng} initialised with a recorded
seed.
\Ensure Estimate $\hat C$, half-width $h_M$, sample variance $S_M^2$.
\State $h \gets T / N$,
$\quad s \gets \sigma \sqrt{h}$,
$\quad D \gets e^{-rT}$
\For{$i = 1, \ldots, M$}
\State $S^{(i)} \gets S_0$
\For{$n = 0, 1, \ldots, N-1$}
\State Draw $\Delta W \sim \mathcal{N}(0, h)$ using \texttt{rng}
\State $S^{(i)} \gets S^{(i)}\,(1 + r h + \sigma\, \Delta W)$
\EndFor
\State $Y_i \gets D \cdot \max(S^{(i)} - K,\, 0)$
\EndFor
\State $\hat C \gets M^{-1} \sum_i Y_i$
\State $S_M^2 \gets (M-1)^{-1} \sum_i (Y_i - \hat C)^2$
\State $h_M \gets z_{\beta/2}\, \sqrt{S_M^{\,2}/M}$
\State \Return $(\hat C,\, h_M,\, S_M^{\,2})$
\end{algorithmic}
\end{algorithm}

The implementation will collapse the inner loop into a single
vectorised pass over all $M$ paths simultaneously, advancing the
state $\hat S_n$ as a length-$M$ array. The pseudocode is written
elementwise for clarity only.

\section{Practical considerations: choosing $h$ and $M$}
\label{sec:practical}

By the MSE decomposition~\eqref{eq:mse-decomp},
\begin{equation}
\E\big[\,(\hat C_{M, h}^{\mathrm{Euler}} - C_{\mathrm{BS}})^2\,\big]
\;\sim\; \frac{C_1}{M} + C_2\, h^2.
\end{equation}
The total computational cost per path is proportional to $N = T/h$
(the inner time loop), so the overall cost of the simulation is
proportional to $M \cdot N = M T / h$. To minimise MSE at a fixed
budget $W = M / h$ (proportional to total cost), set $L(h) = C_1 h /
W \cdot \text{(rearrangement)} + C_2 h^2$, take the derivative with
respect to $h$, and equate to zero. The optimal allocation balances
$\frac{C_1}{M} \sim C_2 h^2$, i.e.\ $M \sim h^{-2}$, equivalent to
\begin{equation}
\label{eq:optimal-allocation}
M^{\,*}(h) \;\propto\; h^{-2},
\quad \text{or equivalently,} \quad
h^{\,*}(M) \;\propto\; M^{-1/2}.
\end{equation}
At the optimum, both terms in the MSE are of the same order, and
the achieved error scales as $W^{-1/3}$ in computational budget $W$,
markedly slower than the unbiased $W^{-1/2}$ rate of Block~1.1.

This degradation is the price of using Euler. The fact is well-known
and is, again, the reason no one would actually use Euler for
European pricing under GBM: the unbiased exact sampler of Block~1.1
delivers $W^{-1/2}$, an order of magnitude better at large budgets.
The Euler scheme is operationally relevant when no closed-form
solution exists (local volatility, stochastic volatility,
jump-diffusion), in which case the comparison is not between Euler
and exact but between Euler and other discretisation schemes.

This analysis is also the foundation for the multilevel Monte Carlo
techniques of Giles (2008), which combine simulations at multiple step
sizes in a way that recovers $W^{-1/2}$ even with biased samplers.
Multilevel MC is beyond the scope of Phase~2 but appears prominently
in the path-dependent literature; we may revisit it in a future
phase.

\section{Bibliographic notes}
\label{sec:biblio}

The empirical convergence study reproduced in this writeup is the
classical demonstration of Higham~\cite{Higham2001}, with adapted
parameters. Glasserman~\cite[Chapter~6]{Glasserman} treats Euler
specifically for finance applications and includes the bias--variance
analysis of Section~\ref{sec:practical}. Kloeden and
Platen~\cite{KloedenPlaten} provide the formal proofs of the
convergence orders cited from the previous theoretical writeup.

\begin{thebibliography}{99}

\bibitem{Glasserman}
P.\ Glasserman,
\emph{Monte Carlo Methods in Financial Engineering},
Springer, 2004.

\bibitem{Higham2001}
D.\ J.\ Higham,
\emph{An algorithmic introduction to numerical simulation of
stochastic differential equations},
SIAM Review, 43(3):525--546, 2001.

\bibitem{KloedenPlaten}
P.\ E.\ Kloeden and E.\ Platen,
\emph{Numerical Solution of Stochastic Differential Equations},
Springer, 1992.

\bibitem{Etheridge}
A.\ Etheridge,
\emph{A Course in Financial Calculus},
Cambridge University Press, 2002.

\end{thebibliography}

\end{document}
